\section{Weakenable families}\label{sec:weakenable}

The localization of Part~VIII is of radius type: a piece that depends on the data in a ball of radius $m$ is bounded by $e^{-\gamma m}$. The polymer expansion of the next stage needs volume type, a fixed small factor for every cube of the dependence region, as in \cite[(116)--(122)]{DimockI}. Ba\l aban and Dimock obtain this from random walks with weakening parameters $s_\Box$, with analyticity for $|s_\Box|\le e^{\kappa}$. This part constructs such weakened walks for all the operators of the step, in $L^2$ norms. Section~\ref{sec:pointwise-weak} explains what is missing for pointwise norms.

\paragraph{Setting.} We keep the conventions of Part~VIII. The background $U$ is real and $(M,\eps)$-regular, or a complex perturbation of such a background as in Proposition~\ref{prop:reg-complex}. Here $\rho$ is a fixed number, chosen in Section~\ref{sec:nested-walks}, independent of $M$. The torus side satisfies $n\ge C(\rho+C_v)M$, so that the sets below do not wrap around. We use the grid points $z_j$, the functions $\chi_j$, $h_j=\chi_j^2$ and the sets $K_j$, $\Lam_j$ of Section~\ref{sec:domains}, and the enlarged domains
\[
\Lam^v_j=\text{the union of the blocks meeting }\{x:|x-z_j|_\infty\le\rho M\}.
\]
For $\rho=9$ these are the domains $\Lam^w_j$ of Section~\ref{sec:random-walk}. We have $|h_j(x)-h_j(x')|\le CM^{-1}|x-x'|$, and every point lies in at most $3^d$ supports $K_j$.

\paragraph{Boxes and polymers.} Fix a \emph{product paving} $\mathcal P=\{I^1_{a_1}\times\dots\times I^d_{a_d}\}$ of $\Tf$. Here each $\{I^\mu_a\}_a$ is a partition of the circle $\R/n\Z$ into intervals of integer length between $M/4$ and $M$. Its elements are called \emph{boxes}; they are rectangular unions of unit blocks. A \emph{polymer} is a $*$-connected union $X$ of boxes: any two boxes of $X$ are joined by a chain of boxes of $X$, consecutive ones sharing at least a corner. Because the paving is a product, two boxes that do not touch are separated in some direction by a whole interval, so they are at distance at least $M/4$. At most $2^d$ boxes lie within distance $r<M/8$ of a point, and each box touches at most $3^d$ boxes. Write $|X|$ for the number of boxes of $X$, and $X^{(r)}$ for the $r$-neighbourhood in lattice units. For each $j$ let $X(j)$ be the union of the boxes that meet $(\Lam^v_j)^{(R_v)}$. Here $R_v\le C_vM$ is a fixed margin, chosen at each level to contain the dependence region of the local propagators. The local operators on a domain $Y$ depend on the background only on $\hat Y$, and those of the fluctuation operator only on $\widehat{Y^+}$, with $Y^+=Y^{(R_0)}$ and the margin $R_0\ge R_0(L,M)$ of Proposition~\ref{prop:local-fluct}(b). We take $R_v\ge R_0+4+2r_{\rm loc}$, where $r_{\rm loc}$ is the largest range of the local operators of the step; then $(\Lam^v_j)^{(R_v)}$ contains these sets for $Y=\Lam^v_j$, together with their $2r_{\rm loc}$-neighbourhoods. Then $X(j)$ is a polymer, and $|X(j)|\le\mathsf N_{\rm b}$, with $\mathsf N_{\rm b}$ depending only on $d$, $\rho$ and $C_v$. (The subscript distinguishes this number of boxes from the matrix size $\mathsf N$ of $G\subset U(\mathsf N)$.)

\paragraph{Monomials.} For $s=(s_\Box)_{\Box\in\mathcal P}\in\mathbb C^{\mathcal P}$ and a function $n:\mathcal P\to\{0,1,2,\dots\}$ with finite support, put $s^n=\prod_\Box s_\Box^{n(\Box)}$, with $\deg s^n=\sum_\Box n(\Box)$ and $\supp s^n=\{\Box:n(\Box)\ge1\}$. For a union of boxes $X$, $s^X$ denotes the monomial with $n=\one_X$. For $\kappa>0$ let $\mathbb D_\kappa=\{s:|s_\Box|\le e^\kappa\text{ for all }\Box\}$. Repeated factors in a monomial are harmless; only the support and the degree enter.

\subsection{Decay norms}

The field spaces are those of Parts~II--VII: scalar fields, gauge fields and $2$-forms on $\Tf$, and unit-lattice fields on $\T$, including fields on the free bonds $\mathfrak F$. All of them decompose over unit blocks, with projections $\one_y$, $y\in\T$. For an operator $F$ between two such spaces and $\gamma\ge0$ put
\[
|F|_\gamma=\max\Big\{\sup_y\sum_{y'}e^{\gamma|y-y'|}\|\one_yF\one_{y'}\|,\ \sup_{y'}\sum_ye^{\gamma|y-y'|}\|\one_yF\one_{y'}\|\Big\}.
\]

\begin{lemma}[Decay norms]\label{lem:decay-norm}
\begin{enumerate}[label=(\alph*)]
\item $\|F\|\le|F|_0\le|F|_\gamma$, and $|FG|_\gamma\le|F|_\gamma|G|_\gamma$.
\item If $F$ has decay $(C,\gamma')$ with $\gamma'>\gamma$, then $|F|_\gamma\le CS_{\gamma'-\gamma}$, with $S_\beta=\sum_ye^{-\beta|y|}$.
\item $|[F,h_j]|_{\gamma/2}\le C(\gamma M)^{-1}|F|_\gamma$ for $0<\gamma\le1$, with $C$ depending only on $d$.
\end{enumerate}
\end{lemma}

\begin{proof}
(a) The first inequality is the Schur test. For the second, use
\[
\|\one_yFG\one_{y''}\|\le\sum_{y'}\|\one_yF\one_{y'}\|\,\|\one_{y'}G\one_{y''}\|,\qquad e^{\gamma|y-y''|}\le e^{\gamma|y-y'|}e^{\gamma|y'-y''|}.
\] (b) is immediate. (c) Write
\[
\one_y[F,h_j]\one_{y'}=\one_yF\one_{y'}\big(h_j-h_j(c_y)\big)-\big(h_j-h_j(c_y)\big)\one_yF\one_{y'} .
\]
On $B(y)$ we have $|h_j-h_j(c_y)|\le CM^{-1}$, and on $B(y')$ we have $|h_j-h_j(c_y)|\le CM^{-1}(1+|y-y'|)$. Hence $\|\one_y[F,h_j]\one_{y'}\|\le CM^{-1}(1+|y-y'|)\|\one_yF\one_{y'}\|$. Finally $(1+t)e^{-\gamma t/2}\le C\gamma^{-1}$.
\end{proof}

On the fine lattice $D$ is not bounded uniformly in $\eta$. As in Parts~II and~III, operators are therefore estimated together with the versions carrying derivatives. For an operator $F$ on fine fields put
\[
|F|^D_\gamma=|F|_\gamma+|DF|_\gamma+|FD^T|_\gamma+|DFD^T|_\gamma ,
\]
where $D$ stands for the covariant derivative appropriate to the spaces involved. For scalar operators followed by $D$ on gauge fields, and so on, the obvious variants are meant.

\subsection{Weakenable families}

\begin{definition}[Weakenable family]\label{def:weak-family}
Let $r\le M/16$. A \emph{weakenable family with margin $r$} is a collection $(X_\omega,n_\omega,O_\omega)_{\omega\in\Omega}$ with the following properties:
\begin{itemize}
\item each $X_\omega$ is a polymer;
\item $n_\omega$ is a function on $\mathcal P$ with $\supp s^{n_\omega}$ equal to the set of boxes of $X_\omega$;
\item each $O_\omega$ is an operator whose kernel is supported in $X_\omega^{(r)}\times X_\omega^{(r)}$, and which depends on the background, and on any other data, only on $X_\omega^{(r)}$.
\end{itemize}
We write $O(s)=\sum_\omega s^{n_\omega}O_\omega$, $\deg\omega=\deg s^{n_\omega}$, and
\[
\|O\|_{\kappa,\gamma}=\sup_{\Box\in\mathcal P}\sum_{\omega:\,\Box\subset X_\omega}e^{\kappa\deg\omega}|O_\omega|_\gamma ,\qquad\|O\|^D_{\kappa,\gamma}=\text{the same with }|O_\omega|^D_\gamma .
\]
We say that $O(s)$ \emph{represents} an operator $O$ if $O(1)=O$, the series converging absolutely.
\end{definition}

The data dependence is part of the definition. It is what makes the weakening also localize the dependence on the background. A family with margin $r$ can always be converted to one with margin $0$, on the level of boxes. Replace $X_\omega$ by $X^\sharp_\omega$, the union of the boxes meeting $X_\omega^{(r)}$, and multiply the monomial by $s^{X^\sharp_\omega\setminus X_\omega}$. Degrees grow by at most a factor $3^d$, and every term then depends on the data, and acts on fields, only inside $X^\sharp_\omega$. We use this conversion whenever dependence inside a union of boxes, not a neighbourhood of it, is needed. This is the case for the factorization over disjoint polymers in the cluster expansion.

\emph{Additive local terms} are the terms of a family that come neither from a walk (Theorem~\ref{thm:weak-inverse}), nor from a product of families (Lemma~\ref{lem:weak-chain}), nor from a local operator applied to a family (Lemma~\ref{lem:weak-local-ops}). In Section~\ref{sec:nested-walks} they are the identities in $R_X=\one_X-P_{\cV,X}$ and $\Pi_X=I-DS_X$ (level~(L2)), the local part $d^2\mathcal N(0)|_X$ of $W_2^X$ (level~(L4)), and the constant $-a$ in $\Delta^{0,X}=\boldsymbol\Gamma^X[W_2^X]-a$ (the repetition of level~(L4)). The constant $r$ of level~(L5) belongs to the local part of the operator there and is not a term of a family. Such a term $T$, a multiple of the identity or a local operator of range at most $r$, is split canonically, box by box. A bond, a plaquette or a unit bond is based at its initial point, and a point at itself. For a box $\Box$ let $\one_\Box$ be the projection onto the field components at the sites based in $\Box$. Then
\[
T=\sum_{\Box\in\mathcal P}T\one_\Box ,
\]
and $T\one_\Box$ is a term with polymer $\Box$ and monomial $s_\Box$; it has margin at most $r$, a bound that does not depend on $M$. Since the boxes partition $\Tf$, this is a decomposition of $T$. It uses only the paving and the base points, so it involves no choice; in particular it commutes with the translations that preserve the paving. The split is used in the $L^2$ norms of this part. A split by indicators does not preserve H\"older tiers, so no tier bound is claimed for an additive term standing alone (Section~\ref{sec:tiers}). The constant $\frac1{1+t}$ in the determinant densities is split by the partition of unity $h_j$ (proof of Corollary~\ref{cor:weak-covariance}), which is canonical as well.

\begin{lemma}[Basic properties]\label{lem:weak-basic}
Let $O$ be a weakenable family with $\|O\|_{\kappa,\gamma}<\infty$.
\begin{enumerate}[label=(\alph*)]
\item For $s\in\mathbb D_\kappa$ the series $O(s)$ converges absolutely, $|O(s)|_\gamma\le2^d\|O\|_{\kappa,\gamma}$, and $s\mapsto O(s)$ is analytic in the interior of $\mathbb D_\kappa$. The same holds with $|\cdot|^D$.
\item \emph{(Localization.)} Let $Z$ be a union of boxes, and let $s_\Box=0$ for all $\Box\not\subset Z$. Then only the terms with $X_\omega\subset Z$ contribute. $O(s)$ is block diagonal with respect to the $*$-components $Z_i$ of $Z$, i.e.\ $\one_{Z_i^{(r)}}O(s)\one_{Z_{i'}^{(r)}}=0$ for $i\ne i'$, and $O(s)$ vanishes off $Z^{(r)}\times Z^{(r)}$. The block on $Z_i$ depends on the data only on $Z_i^{(r)}$, and on the weakening parameters only through $s_\Box$, $\Box\subset Z_i$.
\item \emph{(Covariance.)} If every $O_\omega$ is covariant under real gauge transformations, so is $O(s)$.
\end{enumerate}
\end{lemma}

\begin{proof}
(a) A term contributes to $\one_yO(s)\one_{y'}$ only if $y\in X_\omega^{(r)}$. Then $X_\omega$ contains one of the at most $2^d$ boxes within distance $r$ of $y$. Hence
\[
\sum_{y'}e^{\gamma|y-y'|}\|\one_yO(s)\one_{y'}\|\le\sum_{\Box\text{ near }y}\ \sum_{\omega:\Box\subset X_\omega}|s^{n_\omega}|\,|O_\omega|_\gamma\le2^d\|O\|_{\kappa,\gamma},
\]
and similarly for the column sums. The convergence is uniform on $\mathbb D_\kappa$, so the sum is analytic. (b) If $X_\omega\not\subset Z$, then $s^{n_\omega}=0$. A polymer $X_\omega\subset Z$ lies in one component $Z_i$, and its monomial involves only the boxes of $X_\omega$. Distinct components are at distance at least $M/4>2r$, so $X_\omega^{(r)}$ meets no other $Z_{i'}^{(r)}$. (c) is clear.
\end{proof}

\begin{lemma}[Products]\label{lem:weak-chain}
Let $F_1,\dots,F_n$ be weakenable families with margins at most $r$, and let $0<\delta\le1$. Index the product $F_1\cdots F_n$ by the tuples $(\omega_1,\dots,\omega_n)$ for which $X^{(r)}_{\omega_i}$ and $X^{(r)}_{\omega_{i+1}}$ meet for every $i$. All other products vanish, and this index set is determined by the geometry alone, not by the background. Give each tuple the domain $\bigcup_iX_{\omega_i}$ and the monomial $\prod_is^{n_{\omega_i}}$. This is a weakenable family with margin $r$, and
\[
\|F_1\cdots F_n\|_{\kappa,\gamma}\le n\prod_{i=1}^n\big(C_d\delta^{-2}\|F_i\|_{\kappa+\delta,\gamma}\big).
\]
\end{lemma}

\begin{proof}
If the product of the terms does not vanish, consecutive kernel supports meet. So $X_{\omega_i}^{(r)}\cap X_{\omega_{i+1}}^{(r)}\ne\emptyset$, and the boxes involved are within distance $2r<M/4$. Hence they touch, and the union is a polymer. The monomial has the right support, and the product of the terms has kernel in the union of the supports and depends only on the data there. By Lemma~\ref{lem:decay-norm}(a), $|\prod_iF_{i,\omega_i}|_\gamma\le\prod_i|F_{i,\omega_i}|_\gamma$.

Fix a box $\Box$ and let $i$ be the first index with $\Box\subset X_{\omega_i}$. Given $\omega_i$, sum over $\omega_{i+1}$: it must meet $X_{\omega_i}^{(2r)}$, which is covered by at most $3^d|X_{\omega_i}|$ boxes. So the sum is at most $3^d|X_{\omega_i}|\,\|F_{i+1}\|_{\kappa,\gamma}$. Continue to the right, and in the same way to the left of $i$. Each factor receives at most two factors $|X_{\omega_l}|\le\deg\omega_l$, from its neighbours. Since $\sum_\omega e^{\kappa\deg\omega}(\deg\omega)^2|F_\omega|\le C\delta^{-2}\sum_\omega e^{(\kappa+\delta)\deg\omega}|F_\omega|$, each factor contributes at most $C_d\delta^{-2}\|F_l\|_{\kappa+\delta,\gamma}$. Summing over $i$ gives the factor $n$.
\end{proof}

\begin{lemma}[Local factors and commutators]\label{lem:weak-local-ops}
Let $O$ be a weakenable family with margin $r$.
\begin{enumerate}[label=(\alph*)]
\item Let $\mathcal L$ be a bounded operator of range $r_0$ that depends on the data only within distance $r_0$ of its support. Examples are $\bQ$, $\bQ^T$, $Q'$, $Q'^T$, $D_V$, $\mathfrak F$, $W\in\cW_{\rm loc}$, multiplication by $h_j$ or $m_X$, and the elimination maps $\hat Z^1$, $\hat Z^2$. Then $\mathcal LO$ and $O\mathcal L$, with the same domains and monomials, are weakenable with margin $r+2r_0$, and $\|\mathcal LO\|_{\kappa,\gamma}\le|\mathcal L|_\gamma\,\|O\|_{\kappa,\gamma}$, with $|\mathcal L|_\gamma\le e^{\gamma r_0'}|\mathcal L|_0$, where $r_0'$ is the range in block units. The block Schur norm $|\mathcal L|_0$ is at most $C(d,r_0')\max_{y,y'}\|\one_y\mathcal L\one_{y'}\|$. It can exceed the operator norm $\|\mathcal L\|$, so the latter cannot be used here.
\item $[O,h_j]$, with the same domains and monomials, is weakenable with $\|[O,h_j]\|_{\kappa,\gamma/2}\le C(\gamma M)^{-1}\|O\|_{\kappa,\gamma}$. The same holds with $|\cdot|^D$, up to an additional term $CM^{-1}\|O\|^D_{\kappa,\gamma}$ coming from $[D,h_j]$.
\end{enumerate}
\end{lemma}

\begin{proof}
(a) Apply Lemma~\ref{lem:decay-norm}(a) term by term. The supports grow by the range of $\mathcal L$. (b) Apply Lemma~\ref{lem:decay-norm}(c) term by term. For the derivative versions use $D[O_\omega,h_j]=[DO_\omega,h_j]-[D,h_j]O_\omega$, where $[D,h_j]$ is a multiplication by differences of $h_j$, of size $CM^{-1}$.
\end{proof}
